\section{Study Design Checklist}
\label{sec:Checklist}

This checklist summarizes considerations relevant for any study design using IG 2.0. The checklist is not prescriptive, but rather a collation of the methodological aspects introduced throughout the codebook, serving the development of a project-specific coding protocol. Reflecting on these aspects in the study design phase aims at ensuring high levels of reliability and validity of the coded institutional information. 

Overarching the development of a project-specific codebook adaptation is the question as to whether you intend to develop a generically useful dataset (useful for later adoption using specific techniques or for comparative studies), or whether your coding is purpose-specific (i.e., aimed at addressing specific needs as well as relying on specific analytical techniques). This will inform choices with respect to statement types, encoding heuristics, as well as help establish the relevant IG feature set as well as the scope of encoding, both in terms of number of statements and interpretational scope (see below).

\begin{itemize}
\item Identify applicable statement types (regulative, constitutive, hybrids)
\item Identify level of expressiveness (IG Core, IG Extended, IG Logico), or customized characterizations
\begin{itemize}
\item See \cref{subsec:IgProfiles}
\end{itemize}
\item Decide on the interpretational scope of statement encoding (narrow vs.~wide)
\begin{itemize}
\item Note: This is of specific relevance when coding both constitutive and regulative statements, since it influences the decision heuristics.
\item See \hyperlink{tablink:InterpretationalScopeOfInstitutionalStatements}{here} (\cref{subsec:StatementTypesHeuristics})
\end{itemize}
\item Decide on applicable heuristics for Attributes/Attributes Properties, Objects/Object properties, Constituted Entity/Constituted Entity Properties, and Constituting Properties/Constituting Properties Properties respectively
\begin{itemize}
\item See \hyperlink{tablink:AttributePropertyHeuristics}{here} (\cref{subsubsec:IGCoreRegulative})
\end{itemize}
\item Decide on specific pre-coding/processing steps based on the suggestions listed in \cref{sec:PrecodingSteps}
\begin{itemize}
\item This is particularly relevant where potential reconstruction of institutional statements (conceptual reification) is considered prior to coding.\footnote{For a detailed discussion of conceptual reification, see also Chapters 3 and 5 in \citet{Frantz2022InstitutionalGrammar}.}
\item Decide on aspects such has component scoping during encoding. This includes the decision whether prepositions (`by', `above', `beyond') and articles (`a', `the') are included in the encoding. For instance, if employing modelling techniques, considering such information is counter-productive in the encoding process. An alternative vehicle is the use of semantic annotations to ensure unambiguous labelling of specific component values, especially where you experience divergence in expression (e.g., `EU member states', `member states', `members').
\item Further consider project-specific preparation and coding conventions under consideration of domain-specific and epistemic considerations.
\end{itemize}
\item Decide on applicable taxonomies (see \cref{sec:Taxonomies}), as well as possible extensions and/or additional taxonomies (derived from domain-specific theories or frameworks, for instance).
\end{itemize}

